% Updated draft, 2026-10-27. Rebuilt from the PDF draft (no .tex source was in the repo) and
% filled with results from record/tables/SUMMARY.md and record/tables/*.md (AdilTK/latent-skills,
% branches main and adil) and record/adilet/ (tasadilet-ctrl/latent-skills-freemath).
% Tables in tables/*.tex are generated by make_tables.py from the saved per-question cells.
% Compile with the ACL style files (acl.sty, acl_natbib.bst) in the same folder, e.g. on Overleaf.
\documentclass[11pt]{article}
% preprint shows the authors; use [review] for anonymous submission
\usepackage[preprint]{acl}
\usepackage{times}
\usepackage{latexsym}
\usepackage[T1]{fontenc}
\usepackage[utf8]{inputenc}
\usepackage{microtype}
\usepackage{inconsolata}
\usepackage{graphicx}
\usepackage{booktabs}
\usepackage{amsmath}
\usepackage{amssymb}
\usepackage{multirow}
\usepackage{xcolor}
\usepackage{xurl}  % lets long URLs in the references break

\newcommand{\open}[1]{\textcolor{red}{[OPEN: #1]}}
\newcommand{\TS}{\textsc{ts}}
\newcommand{\ES}{\textsc{es}}
\newcommand{\nTextTests}{147}\newcommand{\nTextSig}{16}\newcommand{\nTextHolmRow}{9}\newcommand{\nTextHolmAll}{4}
\newcommand{\nVecTests}{30}\newcommand{\nVecSig}{3}\newcommand{\nVecHolmRow}{3}\newcommand{\nVecHolmAll}{2}


% \open{Title is a proposal; the draft had none.}
\title{Text Skills and Steering Vectors Give No Measurable Gain on\\Math Benchmarks Once Answer Shortcuts Are Removed}

\author{Adilet Tasbolat$^{*}$ \quad Adil Taibassarov$^{*}$ \quad Rakhat Meiramov$^{*}$ \\
  ISSAI, Nazarbayev University, Astana, Kazakhstan \\
  {\normalfont\small $^{*}$Equal contribution.} \\
  \textit{Preprint, work in progress. October 2026}}

% \input is not allowed inside a tabular before \bottomrule; the primitive is.
\makeatletter\newcommand{\tablerows}[1]{\@@input #1 }\makeatother

\begin{document}
\maketitle

\begin{abstract}
Agent skills are procedures written as text and loaded into the context. We ask whether the same
procedure can instead be delivered as an activation vector, and whether either form transfers
across domains. We compare text skills with contrastive steering vectors on LiveMathematicianBench,
SpreadsheetBench and OfficeQA with Qwen3-8B, Qwen3-14B and Qwen3.6-27B, and we separate the
procedure from the bare presence of a system prompt with a length-matched placebo, both as text and
as an extraction baseline. Two findings reshape the question. First, every large LiveMath gain,
including that of a skill released by an automatic skill optimiser, comes from a benchmark artefact:
an answer option that is correct on all 263 of the 608 items where it appears. Second, once that
artefact is controlled, no skill survives its placebo. We then build a shortcut-free suite of seven
math benchmarks (two multiple-choice sets audited for answer-structure shortcuts and five
free-form ones), learn skills on it with SkillOpt under a pre-registered three-model selection
rule, and test text skills, placebos and steering vectors with paired exact tests. Of
\nTextTests{} text-skill comparisons, \nTextSig{} are significant and all but six are losses; five of the
six gains are on one benchmark and model (HARDMath2, Qwen3.6-27B), where a length-matched placebo
gains more than any skill; of
\nVecTests{} vector comparisons, \nVecSig{} are significant and all are losses. Vectors reproduce a
third of a text skill's effect on the next-token distribution, but the text skill itself changes
no accuracy, so there is nothing for them to carry.
\end{abstract}

\section{Introduction}

Agent skills package a procedure as a Markdown document that is loaded into the context when it is
relevant \citep{anthropic2025skills}. A growing literature evolves such documents automatically from
trajectories and reports gains on reasoning and agentic benchmarks
\citep{alzubi2026evoskill,ni2026trace2skill,yang2026skillopt,tang2026wikiskill}. Two questions are
left open by that work. Every skill is text, so the procedure cannot be separated from the tokens
that carry it; and transfer is reported as a headline gain, without a control for the generic effect
of adding \emph{an} instruction to the prompt, which is known to move accuracy on its own
\citep{webson2022prompts,sclar2024formatting}.

We ask whether a procedure written for one benchmark helps on others, and whether the answer depends
on how it is delivered: as text in the prompt, or as a vector added to the residual stream
\citep{panickssery2024caa}. Our original design was a $3\times3$ matrix of benchmarks by skill
source, with each cell evaluated under a text skill, two steering vectors and a placebo control.
Running it changed the question. On LiveMathematicianBench \citep{he2026livemath}, 263 of the 608
items contain the option ``One of the remaining options is correct, but a stronger result can be
proven'', and it is the correct answer every time. Every large gain we observed on LiveMath, from a
hand-written skill, from SkillOpt's released skill and from our own SkillOpt runs, is a rule for
choosing that option; on the remaining items every skill is matched by a placebo of the same length.
We therefore built a suite of seven math benchmarks without such shortcuts and repeated the study
there, with skills learned by SkillOpt and selected by a rule fixed in advance.

\paragraph{Contributions.}
(1) A text-versus-vector comparison of procedural skills across three benchmarks and three models,
with a placebo-subtracted extraction arm that separates the procedure from prompt presence
(\S\ref{sec:matrix}).
(2) Evidence that every large LiveMath skill gain, including that of a published skill optimiser,
is an answer-format shortcut, and that the remainder does not exceed a placebo
(\S\ref{sec:shortcut}).
(3) A shortcut-free suite of seven math benchmarks with a pre-registered skill-selection protocol,
on which no text skill and no steering vector beats a length-matched placebo
(\S\ref{sec:clean}--\ref{sec:cleanvec}).
(4) Paired exact significance tests with Holm correction throughout, and a list of harness faults
that each would have reversed a conclusion without raising an error (Appendix~\ref{app:faults}).

\section{Related Work}

\paragraph{Skills as text.} Agents have long kept what they learn in natural language rather than in
weights, as verbal reflections \citep{shinn2023reflexion}, extracted insights \citep{zhao2024expel},
or libraries of workflows and executable skills
\citep{wang2024voyager,wang2025awm,wang2025asi,zheng2025skillweaver}. Agent Skills standardised the
artefact as a \texttt{SKILL.md} document that is loaded into context on demand
\citep{anthropic2025skills,xu2026agentskills}. A recent line of work builds such documents
automatically from trajectories: EvoSkill evolves skill folders from failure analysis
\citep{alzubi2026evoskill}, Trace2Skill merges patches proposed by trajectory analysts
\citep{ni2026trace2skill}, SkillOpt treats the document as a trainable parameter behind a validation
gate \citep{yang2026skillopt}, and WikiSkill separates a persistent experience wiki from the active
skills \citep{tang2026wikiskill}. Several of these report that skills transfer to other tasks or
models, but every skill is text, so none can separate a procedure from the tokens that carry it. We
adopt WikiSkill's benchmark set so that our text-skill results are comparable with this literature.

\paragraph{The cost of instructions in context.} Irrelevant content degrades reasoning
\citep{shi2023distracted}, long contexts are used unevenly \citep{liu2024lost}, and accuracy is
sensitive to semantically irrelevant formatting \citep{sclar2024formatting}.
\citet{webson2022prompts} find that models learn about as well from irrelevant or misleading
instructions as from instructive ones. Together these suggest that a skill used outside its own
domain may help or hurt for reasons unrelated to its procedure, and that the mere presence of an
instruction is a confound.

\paragraph{Skills out of context.} Instructions can be compiled into weights by context distillation
\citep{snell2022distilling}, prompt compression \citep{mu2023gist} or hypernetworks that emit LoRA
adapters from text \citep{hu2022lora,charakorn2025t2l}. LatentSkill applies this to agent skills,
mapping skill documents to LoRA adapters \citep{yu2026latentskill}; it is our nearest neighbour, but
it edits parameters, requires a trained hypernetwork, and does not test transfer across domains. Our
embedding skills instead edit the hidden state and need no training.

\paragraph{Activation steering.} Adding a direction to the residual stream steers behaviour without
prompt tokens or weight edits
\citep{subramani2022extracting,turner2023actadd,li2023iti,zou2023repe}. We extract skills by
difference of means in the style of Contrastive Activation Addition \citep{panickssery2024caa}, the
choice supported by analyses of linear representations
\citep{marks2024geometry,park2024linear,im2025steering} and by single-direction accounts of behaviour
\citep{arditi2024refusal}. Closer to procedures, function vectors \citep{todd2024function}, task
vectors \citep{hendel2023taskvectors} and in-context vectors \citep{liu2024icv} show that
demonstrations can be replaced by a vector, and \citet{stolfo2025instruction} steer instruction
following with vectors computed from inputs with and without an instruction. Steering has limits:
prompting outperforms representation steering on AxBench \citep{wu2025axbench}, and steering vectors
are unreliable across inputs \citep{tan2024reliability}.

\paragraph{Procedural skills as vectors.} \citet{jiang2026latent} are closest to our method. They
extract procedural skill directions on synthetic arithmetic tasks, find that composed vectors
slightly outperform text descriptions of composite skills on Qwen3-0.6B, and transfer the directions
to GSM8K and MATH500, where gains shrink with scale. In the terms of our design, they test only the
diagonal of a skill-by-benchmark matrix within one domain, without significance tests or a control
for prompt presence. We test text and embedding skills over benchmarks from different domains, so
that off-diagonal cells measure transfer to domains a skill was not designed for, and we separate the
procedure from prompt presence with a placebo-subtracted extraction arm.
Appendix~\ref{app:related} gives an extended review.

\section{Method}

We ask whether a procedural skill written for one benchmark helps on the others, and whether that
answer depends on how the skill is delivered: as text in the prompt, or as a vector added to the
residual stream. We compare the two deliveries in a matrix of benchmarks by skill source, and we
separate the effect of the procedure from the effect of merely having a system prompt with a placebo
control.

\subsection{Two representations of one skill}

\paragraph{Text skills (\TS).} For each benchmark $j$ a text skill $\TS_j$ is a short Markdown
document describing a procedure for that benchmark's tasks. It is delivered as the system message;
nothing else in the prompt changes. Each \TS{} was written by hand by the authors from the benchmark's task format, without
per-item answers, and is 255--360 words long. The LiveMath skill of the matrix (\texttt{compare\_claims}) went
through three versions, of which we use the third. It was written after inspecting the aggregate
distribution of distractor types chosen by the first two versions on the evaluation items (never
per-item labels), so we declare it test-set-informed and report every LiveMath result on a held-out
subset as well (\S\ref{sec:stats}; history in Appendix~\ref{app:livemathskill}). A later LiveMath
skill, \texttt{livemath\_ts}, was written without access to any split and contains no rule about the
meta option (\S\ref{sec:shortcut}); it is the skill we carry to the shortcut-free suite.

\paragraph{Embedding skills (\ES).} An embedding skill is a single direction in the residual stream,
obtained by Contrastive Activation Addition \citep{panickssery2024caa}, i.e.\ a difference of means
\citep{marks2024geometry,arditi2024refusal}. Let $h_\ell(x)$ be the residual-stream activation at
layer $\ell$ at the final prompt token of input $x$, and let $\mathcal{P}_j$ be a pool of $N=50$
prompts from benchmark $j$. We render every pool prompt exactly as the evaluator renders it (for
LiveMath: question, all five options and the answer cue; for SpreadsheetBench: instruction, workbook
dump and the answer cue), because a vector extracted from bare questions would be built from a
context the model never sees at test time. Given a system message $s$ and a baseline $b$,
\begin{equation}
  v_\ell(s,b) = \frac{1}{N}\sum_{x\in\mathcal{P}_j}\big[h_\ell(s\oplus x) - h_\ell(b\oplus x)\big].
  \label{eq:caa}
\end{equation}
At inference we add the vector at every token position of layer $\ell$, scaled relative to the norm
of the activation it is added to:
\begin{equation}
  h_\ell \leftarrow h_\ell + \gamma\,\lVert h_\ell\rVert\,\frac{v}{\lVert v\rVert}.
  \label{eq:inject}
\end{equation}
Because $v$ is normalised, its magnitude plays no role: two vectors compared at the same $\gamma$ are
pushed equally hard, and they differ only in direction. The embedding condition uses no system
message and adds no prompt tokens.

\subsection{Extraction arms and the placebo control}
\label{sec:arms}

The baseline $b$ in Eq.~\ref{eq:caa} decides what the vector contains, so we treat it as an
experimental variable with three settings. The \emph{placebo} $p$ is a system prompt matched to the
skill in length and register but with no procedural content: writing-style rules (on brevity,
notation and register) and no mathematical or task procedure.

\begin{description}
  \item[\ES{} vs-dummy] $v(\TS_j, p)$: skill minus placebo. Our estimate of what the
  \emph{procedure} contributes.
  \item[\ES{} vs-nothing] $v(\TS_j, \varnothing)$: skill minus no system message. An upper bound that
  also contains whatever a system prompt does by being present.
  \item[Placebo vector] $v(p, \varnothing)$: placebo minus no system message. It contains no skill at
  all, and it measures the generic prompt-presence component that separates the other two arms.
\end{description}

Reporting only the vs-nothing arm would attribute to the skill an effect that a content-free vector
also produces; the placebo vector is what prevents that. We check the placebo two ways. Used as a
text system prompt it must not move accuracy, and placebos with unrelated content (brevity; notation;
tense and voice) must align equally with the skill direction, which shows the shared component is
prompt presence rather than leaked content (Appendix~\ref{app:placebo}). As a further control we
repeat the text, embedding and placebo conditions with the skill placed in the user turn instead of
the system message, so that no condition contains a system prompt at all. A text placebo is also
reported as its own condition everywhere: as \S\ref{sec:shortcut} shows, it is the comparison on
which every apparent skill effect stands or falls.

\subsection{The cross-domain matrix}
\label{sec:matrixdesign}

We use three benchmarks from the agent-skills literature \citep{tang2026wikiskill}:
LiveMathematicianBench \citep{he2026livemath}, SpreadsheetBench \citep{ma2024spreadsheetbench} and
OfficeQA \citep{databricks2025officeqa}. Each has its own skill pair $(\TS_j, \ES_j)$. A cell $(i,j)$
evaluates the skills of benchmark $j$ on benchmark $i$; the diagonal is in-domain and the
off-diagonal cells test transfer. In an off-diagonal cell $\ES_j$ is extracted from benchmark $j$'s
pool and injected into benchmark $i$ unchanged, and $\TS_j$ is used verbatim, domain-specific
vocabulary included. Because the pool comes from a different benchmark than the scored items,
extraction and evaluation sets are disjoint by construction in every transfer cell. On the diagonal
the 50 pool items remain in the scored set, and we also report the disjoint remainder.

Every cell has the same five conditions: no skill, \TS, \ES{} vs-dummy, \ES{} vs-nothing and the
placebo vector. Within a model the three vector conditions share one layer and one $\gamma$, so the
only difference between them is the direction added.

\paragraph{Models.} We use Qwen3-8B, Qwen3-14B \citep{qwen3report} and Qwen3.6-27B \citep{qwen36},
loaded in bf16 with HuggingFace \texttt{transformers}, since serving engines without access to the
residual stream cannot run the embedding conditions. Thinking mode is disabled throughout. Models
are assigned per benchmark by headroom, because a skill worth two points cannot register on a model
near chance or near ceiling. On LiveMath (5-way multiple choice, chance 20\%) both the 8B and the 27B
are well clear of chance. On SpreadsheetBench the 27B is at ceiling (83\% on our agent subset), so
we use the 14B. On OfficeQA the 14B scores 0\% and the 27B 68.6\%, so we use the 27B. OfficeQA
documents have a median length of 215k tokens, which fits neither model's context; the matrix
therefore uses an oracle-context variant with the cited bulletin pages in the prompt (239
questions; 7 of 246 are excluded because their source bulletins failed the manifest checksum).
The agentic OfficeQA setting, in which the model searches the full documents with tools, is not
reported here. A vector lives in one model's residual stream, so
in a transfer cell $\ES_j$ is extracted on the \emph{target} benchmark's model, from benchmark $j$'s
pool and skill text.

\subsection{Choosing the layer and strength}
\label{sec:fidelity}

We choose $(\ell,\gamma)$ for each model and vector without a benchmark score in the loop, using
fidelity, which measures how far steering moves the model's next-token distribution toward the one
the text skill induces:
\begin{equation}
  F = 1 - \frac{\mathrm{KL}(P_{\TS}\,\Vert\,P_{\mathrm{steer}})}{\mathrm{KL}(P_{\TS}\,\Vert\,P_{\mathrm{plain}})},
  \label{eq:fidelity}
\end{equation}
averaged over teacher-forced tokens of a probe set drawn from the benchmark's own items and disjoint
from the extraction pool. $F=1$ means steering reproduces the text skill's distribution, $F=0$ means
it achieves nothing the plain model did not, and $F<0$ means it moves away. The denominator, the
\emph{floor KL}, is itself informative: it is how much the skill text perturbs the model before any
steering, and we report it per benchmark.

Fidelity is computed on forced tokens and so cannot see free generation collapsing into repetition.
Every candidate operating point therefore also passes a \emph{degeneration guard}: an unforced
generation run whose distinct-token ratio must stay close to that of the unsteered model. The guard
rejected configurations that fidelity preferred, including $\gamma=0.5$ on both models and a
four-layer span at $\gamma=0.2$ that reduced the 27B to a single repeated token
(Appendix~\ref{app:degen}). We also check that the selected layer is not at the edge of the searched
grid, and extend the grid when it is. The resulting operating points are layer 34 of 36 for the 8B,
layer 37 for the 14B and layer 62 of 64 for the 27B, all at $\gamma=0.3$ (for the 14B on SpreadsheetBench, layers 37--38 of 40 with $\gamma$ between
0.3 and 0.5). In transfer cells and on the
shortcut-free suite we keep these operating points fixed rather than re-selecting them, so no
information about the target benchmark enters, and report fidelity on a target probe as a diagnostic
only.

\subsection{Decoding and scoring}

Decoding is greedy. New-token budgets are 2048 (8B), 3072 (14B) and 4096 (27B) on LiveMath, and 1536
for SpreadsheetBench. The SpreadsheetBench matrix cells use a single-turn variant in which the model
writes one script; an agent variant with up to 20 turns gives the same baseline on the 14B (20.0 vs
20.1) and a same-direction text-skill effect (+4.0 vs +2.5), and the steered arms are run
single-turn because the hooked server is about 500 times slower per token. With greedy decoding a
second run would be identical, so each cell is one run and all reported uncertainty is sampling
error over items; we have no run-to-run variance estimate.

\paragraph{Answer forcing.} Some replies on the direct-answer benchmarks spend the whole budget
deliberating and never commit to an answer. Scoring them as wrong would bias the comparison, because
truncation tracks verbosity and skills change verbosity. We therefore continue any non-committal
reply with an ``Answer:'' cue for a few tokens. This is done identically in every condition, and
inside the steering context in the embedding conditions, so that the continuation is steered like
the reply it continues. We report the forced rate per cell; on LiveMath the parse rate is then 1.0
in every cell.

\paragraph{Benchmark protocols.} LiveMath options are shuffled deterministically per question, since
the released data places the correct answer first in every item. For SpreadsheetBench the model
writes an \texttt{openpyxl} script that is executed and graded against the gold workbook over the
task's answer range. We depart from the official protocol in two ways, identically across
conditions: we ask for computed values rather than formulas, since no recalculation engine is
available, and we state the answer range in the prompt, since that is exactly the region graded.
The grader scores 100\% on gold files and 0.3\% for an agent that does nothing. The single-turn
variant uses the 398 tasks of the verified-400 subset that have a gold file, and the agent variant
200 of them; the verified subset scores one test case per task, so our numbers are not comparable
with published full-benchmark results. Further protocol details and the harness faults we found and
fixed are in Appendix~\ref{app:method}.

\subsection{Statistics}
\label{sec:stats}

All conditions in a cell are scored on the same items, so we compare each condition with the
no-skill baseline, and \TS{} with each \ES{} arm, using the exact McNemar test on discordant pairs
\citep{mcnemar1947,dietterich1998}. We correct with Holm's procedure \citep{holm1979} within each
(benchmark, model) family, whose comparisons share a baseline and answer one question. Because a
reader will also ask whether a finding survives correction across the whole study, we report that as
well for every significant result. Alongside accuracy we report items fixed and broken, the
held-out subset for diagonal cells, the forced rate, and fidelity. Fidelity and accuracy are
reported separately: a vector can reproduce a text skill faithfully and still leave accuracy
unchanged when the text skill itself does not help.

\subsection{Benchmark shortcuts and the shortcut-free suite}
\label{sec:suite}

\paragraph{Two LiveMath shortcuts.} LiveMath has two answer-structure shortcuts that require no
mathematics. The \emph{meta option}, ``One of the remaining options is correct, but a stronger
result can be proven'', appears in 263 of the 608 items and is correct in all 263. The
\emph{centroid}, the option most similar to the others, is correct on 68\% of items with no model at
all, because distractors are usually one-detail edits of the correct statement. We therefore report
LiveMath on two subsets: META (the 263 items with the meta option) and OTHER (the remaining 345),
and on OTHER we split again by whether the centroid is right. Only OTHER measures mathematics.

\paragraph{Benchmarks.} We assembled seven math benchmarks without such shortcuts.
Two are multiple-choice and were audited with no model in the loop for every option-structure
heuristic we could think of (position, length, centroid, meta options): the TheoremQA and SciBench
math items of MMLU-Pro \citep{wang2024mmlupro} (10 options; centroid 11.6\% against 10\% chance) and
the hard math items of SuperGPQA \citep{supergpqa2025} (centroid 18.0\%, so reported split by
centroid as well). Five are free-form, with no options at all: BeyondAIME \citep{beyondaime2025},
HARDMath2 \citep{roggeveen2025hardmath2}, RealMath \citep{zhang2025realmath} (the 387 of 633
arXiv items whose answer is a single expression), IMO-AnswerBench \citep{luong2025imobench} (the 345
single-answer items) and ArxivMath \citep{balunovic2025matharena}. Each is split by a deterministic
hash of the item id into train, dev and test (Table~\ref{tab:suite}). Free-form answers are scored
with Math-Verify \citep{mathverify} with exact rational tolerance; Appendix~\ref{app:faults} lists the
scorer faults we fixed.

\begin{table}[t]
\centering\small
\resizebox{\linewidth}{!}{%
\begin{tabular}{@{}lrrr@{}}
\toprule
Benchmark & Items & Train/dev/test & Answers \\
\midrule
MMLU-Pro math & 505 & 40/100/365 & 10 options \\
SuperGPQA hard math & 500 & 40/100/360 & up to 10 \\
BeyondAIME & 100 & 30/30/40 & integers \\
HARDMath2 & 211 & 40/50/121 & formulas \\
RealMath & 387 & 40/60/287 & expressions \\
IMO-AnswerBench & 345 & 40/100/205 & expressions \\
ArxivMath & 103 & 30/30/43 & expressions \\
\bottomrule
\end{tabular}}
\caption{The shortcut-free suite. Splits are a hash of the item id; ids are released, data is not.}
\label{tab:suite}
\end{table}

\paragraph{Learning skills.} We use SkillOpt \citep{yang2026skillopt} as a draft generator, with
Qwen3.6-27B as the optimiser model and Qwen3-8B as the target (Qwen3-14B on ArxivMath, where the 8B
is at the floor). SkillOpt measures progress against its own seed skill, never against no skill, and
its validation splits are small, so we do not use its choice. A selection rule fixed before any test
run picks among SkillOpt's seed and every accepted draft: the draft with the highest mean dev gain
over no skill, averaged over all three models; on multiple-choice benchmarks a draft is disqualified
if its gain on centroid-wrong items is more than 3 points below its gain on centroid-right items and
not positive; ties go to the shorter skill. The chosen draft is then evaluated once on test.

\paragraph{Conditions.} Every model and benchmark is evaluated with no skill, the length-matched
placebo, a hand-written skill (\texttt{strength\_audit}) with and without its meta-option rule,
\texttt{livemath\_ts}, its own length-matched placebo, \texttt{livemath\_ts} with the meta-option
rule added, and the chosen SkillOpt draft. Text conditions are served with vLLM \citep{kwon2023vllm}
with the same prompt, forcing and budgets as above.

\paragraph{Vector arm.} On the two multiple-choice benchmarks we extract three vectors from
\texttt{livemath\_ts} for each model: B1 (skill minus its placebo, the vs-dummy arm), B2 (skill minus
nothing) and CONTROL (placebo minus nothing), from the benchmark's 40 train items rendered as the
evaluator renders them. The operating points of \S\ref{sec:fidelity} are kept fixed. Since vLLM
cannot steer and its numerics differ from HuggingFace's, every cell of the vector arm, text cells
included, is rerun on the HuggingFace engine.

\section{Results}

\subsection{The matrix}
\label{sec:matrix}

\paragraph{In-domain skills.} Table~\ref{tab:matrix} gives the diagonal. No text skill is
significant: $-0.2$, $+2.3$ and $+1.6$ points on LiveMath for the 8B, 14B and 27B, $+2.5$ on
SpreadsheetBench (14B, single turn; $+4.0$ with the agent harness) and $+1.7$ on OfficeQA (27B). Five of six text cells are positive, as
expected for roughly three under the null, but none reaches $p<0.05$. Among the vectors, only the
uncontrolled vs-nothing arm on the 27B is significant ($+6.25$, $p=0.003$), and it survives Holm
correction within its row. On the 558 LiveMath items outside the extraction pool the same cell is $+6.6$
(Appendix~\ref{app:pool}), so the overlap confers no advantage.

\begin{table*}[t]
\centering\small
\resizebox{\linewidth}{!}{%
\begin{tabular}{@{}llr r rrrr@{}}
\toprule
Benchmark & Model & $n$ & No skill & \TS & \ES{} vs-dummy & \ES{} vs-nothing & Placebo vector \\
\midrule
LiveMath & Qwen3-8B & 608 & 26.2 & $-$0.2 & $-$0.3 & +1.0 & -- \\
LiveMath & Qwen3-14B & 608 & 25.3 & +2.3 & $-$0.2 & +0.8 & -- \\
LiveMath & Qwen3.6-27B & 608 & 33.1 & +1.6 & +2.1 & \textbf{+6.2$^{*}$} & +3.6 \\
LiveMath, skill in user turn & Qwen3.6-27B & 608 & 33.1 & +2.3 & -- & +2.5 & +4.1$^{*}$ \\
SpreadsheetBench (single turn) & Qwen3-14B & 398 & 20.1 & +2.5 & +0.3 & +2.3 & -- \\
OfficeQA (oracle pages) & Qwen3.6-27B & 239 & 68.6 & +1.7 & -- & -- & -- \\
\bottomrule
\end{tabular}}
\caption{Diagonal of the matrix: change in accuracy (points) against the row's no-skill baseline.
$^{*}$ $p<0.05$, exact McNemar; bold = survives Holm correction within the row; -- = not run. LiveMath here uses
the full 608 items, which mix the META and OTHER subsets (\S\ref{sec:shortcut}).}
\label{tab:matrix}
\end{table*}

\paragraph{Fidelity does not predict effect.} The vs-dummy vectors reproduce the text skill's
distribution to $F=0.43$, $0.64$ and $0.23$ on the 8B, 14B and 27B. The 14B vector, the most faithful
in the study, changes accuracy by $-0.2$ while its text skill gives $+2.3$; its least faithful
vector ($F=0.19$) gives $+0.8$. The floor KL is 0.82--0.89 on LiveMath and 0.19--0.22 on
SpreadsheetBench: the same kind of skill text perturbs the model about four times less in a long
agentic context.

\paragraph{Cross-domain transfer.} Table~\ref{tab:transfer} carries the LiveMath skill
\texttt{livemath\_ts} to OfficeQA and SpreadsheetBench, as text and as both vectors (extracted on a
neutral pool). All changes are within $\pm3$ points; the one significant cell is the 4B text skill on
SpreadsheetBench ($-3.8$). The math skill neither helps nor seriously harms unrelated tasks, and the
vectors are no better than the text. The off-diagonal cells with the SpreadsheetBench and OfficeQA skills as the source, and OfficeQA for
the 14B and 27B, were not run and are left to future work.

\begin{table}[t]
\centering\small
\setlength{\tabcolsep}{4pt}
\resizebox{\linewidth}{!}{%
\begin{tabular}{@{}llrrrr@{}}
\toprule
Target & Model & Base & \TS & vs-dummy & vs-nothing \\
\midrule
OfficeQA & 4B & 23.1 & $-$2.1 & +2.9 & $-$1.7 \\
OfficeQA & 8B & 26.5 & +2.9 & +0.4 & $-$0.8 \\
SpreadsheetB. & 4B & 12.3 & $-$3.8$^{*}$ & $-$3.0 & $-$3.3 \\
SpreadsheetB. & 8B & 15.3 & $-$0.8 & $-$1.0 & +1.5 \\
SpreadsheetB. & 14B & 21.6 & $-$1.5 & $-$2.3 & +0.3 \\
SpreadsheetB. & 27B & 54.0 & +0.5 & +0.5 & +1.3 \\
\bottomrule
\end{tabular}}
\caption{Transfer of the LiveMath skill \texttt{livemath\_ts}. OfficeQA $n=238$, SpreadsheetBench
$n=398$ (single turn). The SpreadsheetBench baselines come from a later harness version than
Table~\ref{tab:matrix}, which is why the 14B baseline is 21.6 here and 20.1 there.}
\label{tab:transfer}
\end{table}

\paragraph{Procedure versus prompt presence.} The placebo vector, which contains no skill, gives
$+3.6$ on LiveMath 27B with the system-message contrast and $+4.1$ ($p=0.04$) with the user-turn
contrast; neither survives correction, but they agree. The vs-dummy arm, which subtracts the
placebo, gives $+2.1$ ($p=0.30$). Most of the vs-nothing gain is therefore prompt presence, not the
procedure: the uncontrolled vector carries most or all of the text skill's effect, and subtracting a
matched placebo removes it. The same holds on SpreadsheetBench ($+2.3$ vs-nothing, $+0.3$ vs-dummy).

\subsection{LiveMath gains are the meta option}
\label{sec:shortcut}

Table~\ref{tab:meta} splits LiveMath by the meta option for every text skill we evaluated, on the 556
items that no SkillOpt run trained on. Every skill containing a meta-option rule gains on META and
loses on OTHER. SkillOpt's released skill, optimised with GPT-5.5, takes META from 3.4\% to 75.4\%
on the 14B and from 16.8\% to 69.0\% on the 27B while OTHER moves by at most 2.5 points; with its two
meta bullets removed, almost the whole gain disappears. Adding those two bullets to
\texttt{livemath\_ts} costs it 4.7 to 5.3 points on OTHER. Our own SkillOpt runs with a Qwen3.6-27B
optimiser learned the same rule: the 8B-trained skill takes META to 68--99\% on every model size
while OTHER stays flat (Appendix~\ref{app:skillopt}).

\begin{table*}[t]
\centering\small
\resizebox{\linewidth}{!}{%
\begin{tabular}{@{}lll rrrrr@{}}
\toprule
Model & Subset & No skill & \texttt{strength\_audit} & SkillOpt (released) & \texttt{livemath\_ts} + meta rule & \texttt{livemath\_ts} & Placebo \\
\midrule
Qwen3-8B & OTHER (324) & 42.3 & $-$4.0 & $-$2.5 & +2.2 & +6.8$^{*}$ & +8.3$^{*}$ \\
 & META (232) & 3.4 & +9.1$^{*}$ & +5.2$^{*}$ & +0.9 & $-$2.6 & $-$2.6 \\
Qwen3-14B & OTHER & 43.5 & $-$6.8$^{*}$ & $-$2.2 & +0.3 & +5.6 & +5.6$^{*}$ \\
 & META & 3.4 & +21.1$^{*}$ & +72.0$^{*}$ & +22.8$^{*}$ & $-$1.7 & +0.4 \\
Qwen3.6-27B & OTHER & 52.5 & $-$4.9 & +0.9 & $-$2.5 & +2.5 & $-$0.3 \\
 & META & 16.8 & +39.7$^{*}$ & +52.2$^{*}$ & +24.1$^{*}$ & $-$6.5$^{*}$ & $-$3.9 \\
\bottomrule
\end{tabular}}
\caption{LiveMath by the meta option (556 items: the 608 minus SkillOpt's 52 train and validation
items). Change in points against no skill; $^{*}$ $p<0.05$. META: the option ``One of the remaining
options is correct, but a stronger result can be proven'', correct whenever present. The placebo
is length-matched to \texttt{livemath\_ts}.}
\label{tab:meta}
\end{table*}

\paragraph{Without the shortcut, the placebo matches the skill.} \texttt{livemath\_ts} is the only
skill positive on OTHER for all three models ($+6.8^{*}$, $+5.6$, $+2.5$), but its placebo, which
contains no procedure, does as well on the 8B and 14B ($+8.3^{*}$, $+5.6^{*}$). Pooled over the three
models, \texttt{livemath\_ts} beats no skill on OTHER by 4.0 points ($p=0.02$), but a length-matched
placebo by only 2.8 ($p=0.10$). The steering vectors of \texttt{livemath\_ts} keep less: none is significant on OTHER,
and the only significant vector cell on the 27B is on META ($+6.5^{*}$, skill minus nothing), a
nudge toward the shortcut rather than mathematics.

\paragraph{Removing the shortcut from training.} SkillOpt skills trained with the meta option kept
out (by dropping META items, filtering edits that mention it, or deleting the rule afterwards) do
not survive the placebo either. The one significant result, the 8B-trained skill with its meta rule
deleted, gives $+5.9^{*}$ on 27B OTHER against $+0.6$ for its placebo; but on the 8B it was trained
on it gives $+2.2$ (placebo $+3.4$), and on a second split of OTHER items the skill and the placebo
tie at $+3.6$. Skills trained on the 27B itself are at most $+2.2$ on 27B OTHER, and the latest
variant (v2) is $-2.6$ against a placebo at $+7.8^{*}$ on the 153 of 205 OTHER items finished so far (the v2 run is incomplete). Appendix~\ref{app:skillopt} gives the tables.

\subsection{Text skills on the shortcut-free suite}
\label{sec:clean}

Table~\ref{tab:clean} gives every text condition on the seven shortcut-free benchmarks.
\textbf{No skill beats its placebo.} Of \nTextTests{} comparisons, \nTextSig{} are significant
at $p<0.05$ (about seven are expected by chance) and eighteen of them are losses. \nTextHolmRow{}
survive Holm correction within their row and \nTextHolmAll{} across the whole table; those four are
all \texttt{livemath\_ts} with and without the meta rule on Qwen3-8B, on RealMath and IMO-AnswerBench,
and they are format failures rather than reasoning losses (below). Six comparisons are gains. One,
the SkillOpt draft on MMLU-Pro with the 14B ($+5.5$), survives correction within its row, did not
appear on dev, and reverses on the 27B with the same text ($-2.5^{*}$). The other five are on
HARDMath2 with the 27B, where every system prompt helps and the length-matched placebo helps most
($+14.9$, against $+8.3$ to $+13.2$ for the skills; below).

\begin{table*}[t]
\centering\small
\setlength{\tabcolsep}{4pt}
\resizebox{\linewidth}{!}{%
\begin{tabular}{@{}llr r rrr rrr r@{}}
\toprule
 & & & & & \multicolumn{2}{c}{\texttt{strength\_audit}} & \multicolumn{3}{c}{\texttt{livemath\_ts}} & SkillOpt \\
\cmidrule(lr){6-7}\cmidrule(lr){8-10}
Benchmark & Model & $n$ & No skill & Placebo & full & $-$meta & placebo & skill & +meta & chosen \\
\midrule
\tablerows{tables/clean_text.tex}
\bottomrule
\end{tabular}}
\caption{Text skills on the shortcut-free suite (test splits): accuracy of no skill, and change in
points for every other condition. $^{*}$ $p<0.05$, exact McNemar, uncorrected; bold = survives Holm
correction within the row. $^{\dagger}$ = at least 20\% of replies answer a free-form question with
an option letter that does not exist. MCQ rows above the rule, free-form rows below. Greedy,
thinking off, answer forcing, 2048 tokens (4096 for the 27B).}
\label{tab:clean}
\end{table*}

\paragraph{The meta rule is inert without a meta option.} \texttt{livemath\_ts} with and without the
meta-option rule differ significantly once in 21 paired comparisons (HARDMath2 27B, $+8.3$ with the
rule, $p=0.013$; next smallest $p=0.07$), on a benchmark that has no options at all. The same rule
took the 14B from 3.4\% to 26.3\% on LiveMath META. Here there is nothing for it to act on.

\paragraph{A multiple-choice skill makes small models invent options.} \texttt{livemath\_ts} says to
compare the options. On free-form questions, with no options, Qwen3-8B often derives the correct
value and then writes ``Among the options, Option A \ldots'' and answers \verb|\boxed{A}|: 186 of 287
replies on RealMath (193 with the meta rule), 59 of 205 on IMO-AnswerBench, 26 of 43 on ArxivMath and
26 of 121 on HARDMath2; the 14B does so on 64 of 287 RealMath items, and the 27B never on more than 6.
The large free-form losses of \texttt{livemath\_ts} are this format failure. Reply length is
unchanged (2586 against 2664 characters on RealMath 8B), so it is not a budget effect.

\paragraph{Placebos hurt by shortening replies.} The length-matched placebos are the most frequent
significant losses (MMLU-Pro 8B $-7.7$ and $-7.1$, SuperGPQA 8B $-6.9$ and $-9.4$, IMO-AnswerBench
8B $-6.3$ and 14B $-6.3$). Their ``be concise'' style rules halve reply length (2586 to 531 characters
on RealMath 8B), and these problems need the room. This is also why ``\texttt{livemath\_ts} beats its
placebo'' (MMLU-Pro 8B, $+54/-24$ discordant items) is not a skill effect: the placebo is harmful,
and \texttt{livemath\_ts} does not beat no skill on any benchmark except HARDMath2 27B.
The 27B on HARDMath2 is the mirror image. With no system prompt its replies are short (median 3{,}563
characters at a 4096-token budget); every system prompt, placebos included, more than doubles them
(8{,}314 to 9{,}437) and raises accuracy, and the placebo raises it most ($+14.9$). Length again
explains the effect; the procedure in the skill adds nothing on top of it.

\paragraph{Optimised drafts do no better than a placebo.} With dev scored by Math-Verify, the
pre-registered rule keeps SkillOpt's four-line seed (or the seed plus empty tags) on four benchmarks
and picks a longer, optimised draft on HARDMath2, RealMath and IMO-AnswerBench. (Scored by string
match only, as first run, it chose the seed on all seven; Appendix~\ref{app:faults}.) The optimised
drafts beat no skill only on HARDMath2 with the 27B ($+13.2$), where the placebo gains more
($+14.9$); on IMO-AnswerBench with the 8B the draft loses $6.3$ points, which survives correction
within its row. On the other seven cells they change accuracy by $-3.1$ to $+2.4$, none significant.
SkillOpt's own validation gains did not reach test, and the optimised drafts contain benchmark
priors and memorised training items rather than general method (Appendix~\ref{app:audit}).

\subsection{Steering vectors on the shortcut-free suite}
\label{sec:cleanvec}

Table~\ref{tab:cleanvec} gives the vector arm on the two multiple-choice benchmarks, with every
cell on the same engine. \textbf{No vector beats no skill.} Of \nVecTests{} comparisons, \nVecSig{}
are significant and all are losses: the text placebo on the 8B (MMLU-Pro $-8.2$, SuperGPQA $-8.3$)
and B1 on SuperGPQA 8B ($-10.3$), which is also below CONTROL ($p=0.04$) and below the text skill
($p<0.001$). The text skill itself changes nothing ($-0.6$ to $+1.9$, none significant), so there is
no gain for a vector to carry, and the vectors carry none.

\begin{table*}[t]
\centering\small
\resizebox{\linewidth}{!}{%
\begin{tabular}{@{}llr r rr rrr@{}}
\toprule
Benchmark & Model & $n$ & No skill & \texttt{livemath\_ts} text & Placebo text & B1 (vs-dummy) & B2 (vs-nothing) & CONTROL \\
\midrule
\tablerows{tables/clean_vector.tex}
\bottomrule
\end{tabular}}
\caption{\texttt{livemath\_ts} as text and as three steering vectors (HuggingFace engine, test
splits). Change in points against no skill; $^{*}$ $p<0.05$; bold = survives Holm within the row.
Vectors: difference of means at the last prompt token over the benchmark's 40 train items;
one layer (8B L34, 14B L37, 27B L62), $\gamma=0.3$, fixed in advance.}
\label{tab:cleanvec}
\end{table*}

\paragraph{Fidelity.} At the fixed operating point the 8B and 14B vectors reach $F=0.23$--$0.42$, and
the 8B CONTROL is $\approx0$ (Table~\ref{tab:fid}): the vectors carry a third of the text skill's
effect on the next-token distribution and the placebo direction carries none. On the 27B at layer 62,
the LiveMath operating point, B1 and B2 reach only $F=0.06$--$0.10$, below CONTROL's $0.20$; that layer
does not carry this skill, so the 27B vector cells are weak evidence.

\begin{table}[t]
\centering\small
\resizebox{\linewidth}{!}{%
\begin{tabular}{@{}llrrrr@{}}
\toprule
Benchmark & Model & $\ell$ & B1 & B2 & CONTROL \\
\midrule
\tablerows{tables/clean_fidelity.tex}
\bottomrule
\end{tabular}}
\caption{Fidelity (Eq.~\ref{eq:fidelity}) at $\gamma=0.3$ on 16 dev items; each vector is measured
from the base it was extracted against.}
\label{tab:fid}
\end{table}

\section{Conclusion}

We set out to compare a procedural skill delivered as text with the same skill delivered as a
steering vector, in and out of domain. Three results answer that question for the settings we
tested. On LiveMath, the large skill gains reported for hand-written and optimised skills alike are
an answer-format shortcut, and without it the remaining gain is matched by a placebo. On seven math
benchmarks without such shortcuts, neither a text skill nor a steering vector improves on the
length-matched placebo for Qwen3-8B, Qwen3-14B or Qwen3.6-27B: across \nTextTests{} text and
\nVecTests{} vector comparisons, most significant effects are losses, and the few gains are matched
or exceeded by the placebo. And what skills and placebos change
is mostly how much the model writes, not how well it reasons. Steering vectors reproduce part of a
text skill's distributional effect, so the method works as intended; there is simply no behavioural
effect to transfer. Our recommendations for skill evaluation follow: audit the benchmark for
answer-structure shortcuts before optimising against it, report a length-matched placebo alongside
every skill, compare against no skill rather than against an optimiser's own seed, and select
skills on dev across models before touching test.

\section*{Limitations}

\textbf{Models.} All models are from the Qwen3 family, and thinking mode is off throughout; a
procedural skill may need reasoning tokens to act on, and other families may respond differently.
\textbf{Optimiser.} Our SkillOpt runs use Qwen3.6-27B as the optimiser; the SkillOpt paper used
GPT-5.5, whose released LiveMath skill we evaluate but did not rerun on the shortcut-free suite.
\textbf{Benchmarks.} The matrix covers three benchmarks, and several of its off-diagonal and
embedding cells were not completed (Table~\ref{tab:matrix}). BeyondAIME ($n=40$) and ArxivMath
($n=43$) can only detect large effects, and the 8B and 14B score below 10\% on HARDMath2 and
ArxivMath with thinking off. \textbf{Vectors.} We test single-direction CAA vectors at one layer and one strength, fixed in
advance; the 27B operating point does not carry \texttt{livemath\_ts}, and per-paragraph vectors or
learned interventions \citep{wu2024reft,wu2025reps} might fare differently. Transfer cells require
re-extraction per model. \textbf{Skills.} The matrix's LiveMath skill is test-set-informed (mitigated
by the held-out subset), and every skill we evaluate is one text; a different skill could behave
differently. \textbf{Runs.} Each cell is one greedy run, so uncertainty is sampling error over items
only.

\bibliography{custom}

\appendix

\section{Extended Related Work}
\label{app:related}

This appendix expands Section 2. We distinguish three tiers of skill representation
(Table~\ref{tab:tiers}); our experiments compare tier 1 with tier 3, and tier 2 is reviewed for
positioning.

\begin{table*}[t]
\centering\small
\begin{tabular}{@{}p{1.9cm}p{3.4cm}p{4.4cm}p{4.9cm}@{}}
\toprule
Tier & Where the skill lives & Inference cost & Representative work \\
\midrule
1. Text & Input tokens in context & Prompt tokens; skill vocabulary in context & Agent Skills, EvoSkill, SkillOpt, WikiSkill \\
2. Weights & Adapter $\Delta W$ & No prompt tokens; one adapter per skill & Text-to-LoRA, LatentSkill \\
3. Activations & Vector $v$ added to $h_\ell$ & No prompt tokens, no weight edits & CAA, function vectors, latent skill steering \\
\bottomrule
\end{tabular}
\caption{Three tiers of skill representation.}
\label{tab:tiers}
\end{table*}

\subsection{Skills as text}

\paragraph{From reflection to skill libraries} The idea that an agent should keep what it learns in
natural language, rather than in its weights, predates the term ``skill''. \emph{Reflexion}
\citep{shinn2023reflexion} has an agent write verbal reflections on its failures into an episodic
memory that conditions later attempts, improving performance without gradient updates.
\emph{ExpeL} \citep{zhao2024expel} generalises this across tasks: the agent gathers experiences,
extracts natural-language insights from successes and failures, and recalls both insights and
trajectories at test time.

A second line stores procedures, not lessons. \emph{Voyager} \citep{wang2024voyager} builds a
growing library of executable code skills while exploring Minecraft, retrieving them by description
and composing them into more complex behaviours. \emph{Agent Workflow Memory} \citep{wang2025awm}
induces reusable workflows from past web-navigation trajectories and adds them to the agent's
context. \emph{Agent Skill Induction} \citep{wang2025asi} represents skills as programs that are
verified by re-execution before being admitted, and \emph{SkillWeaver} \citep{zheng2025skillweaver}
lets web agents discover and hone skills as reusable APIs. Across this line of work, the skill is a
textual or programmatic artefact and its effect is mediated by the context window.

\paragraph{The skill document as a format} Agent Skills, introduced by Anthropic in October 2025
\citep{anthropic2025skills}, standardised the artefact: a skill is a folder of instructions, scripts
and resources organised around a \texttt{SKILL.md} file, and the agent sees only each skill's short
metadata until it decides a skill is relevant, at which point it loads the full instructions and, if
needed, further resources. This \emph{progressive disclosure} is an explicit acknowledgement that
skill text is costly in context. A recent survey \citep{xu2026agentskills} reviews the architecture,
acquisition and security of such skills, reporting that a substantial fraction of community skills
contain vulnerabilities. It treats skills throughout as external packages loaded into context and
does not discuss weight- or activation-space representations.

\paragraph{Automatic skill evolution} A cluster of 2026 papers automates the construction of skill
documents from trajectories. They share an outer loop (run, analyse, edit, validate) and differ in
how they search the space of documents.

\emph{EvoSkill} \citep{alzubi2026evoskill} runs an evolutionary loop over a skill repository. An
executor agent attempts tasks with the current skills; a proposer agent analyses failures and
proposes new skills or edits; a skill-builder agent turns proposals into skill folders containing
trigger metadata, Markdown instructions and optional helper scripts. Candidates are retained on a
fixed-capacity Pareto frontier while the underlying model is held fixed. It reports gains on
OfficeQA and SealQA and zero-shot transfer of a skill evolved on SealQA to BrowseComp.

\emph{Trace2Skill} \citep{ni2026trace2skill} distils a batch of trajectories into a single skill
directory. An error analyst and a success analyst run in parallel over failed and successful
trajectories to propose localised patches, and a hierarchical merge step deduplicates them, resolves
conflicts and generalises across trajectories. It evaluates on spreadsheet tasks (including
SpreadsheetBench), AIME, DocVQA and office-document tasks, and reports that skills evolved with one
model transfer to larger models and other families without parameter updates.

\emph{SkillOpt} \citep{yang2026skillopt} makes the optimisation analogy explicit. A skill document is
trainable external state for a frozen model; an optimiser model proposes bounded add, delete and
replace edits; an ``edit budget'' with a decay schedule plays the role of a learning rate;
candidates are accepted only when they improve a held-out validation split; rejected edits are kept
as negative feedback; and an epoch-level slow update acts like momentum. Its evaluation is the
broadest in this group: SearchQA, SpreadsheetBench, OfficeQA, DocVQA, LiveMathematicianBench and
ALFWorld, across seven models and three execution harnesses. As \S\ref{sec:shortcut} shows, its
LiveMath gain is the meta-option shortcut, which its optimiser found and encoded.

\emph{WikiSkill} \citep{tang2026wikiskill} separates what the agent has learned from what it is
told. A three-layer workspace holds immutable raw traces, a persistent wiki of patterns and
evolution logs, and the active skills injected into the inference agent. A wiki maintainer
consolidates traces; a ReAct-style proposer consults the wiki to propose skill edits; a gate accepts
or rolls back each proposal while the wiki persists. WikiSkill evaluates on LiveMathematicianBench,
SealQA, SpreadsheetBench, OfficeQA and ALFWorld against Trace2Skill, EvoSkill and SkillOpt, on five
models including Qwen3.6-27B. Its benchmark set is the one we adopt.

These systems inherit from the prompt-optimisation literature, in which an LLM proposes prompts
given a history of candidates and scores \citep{yang2024opro}, natural-language feedback is
propagated through a computation graph as a textual ``gradient'' \citep{yuksekgonul2024textgrad}, or
prompts evolve by reflection with Pareto-based candidate selection \citep{agrawal2025gepa}. What is
new in the skills setting is the artefact (a persistent, reusable procedure rather than a task
prompt) and the emphasis on reuse across tasks.

\paragraph{What the text-skill literature leaves open} Three questions are not addressed by this
literature. First, no paper varies the \emph{representation} of the skill: every skill is text, so
the procedure cannot be separated from the tokens that carry it. Second, transfer is reported as a
headline gain, and no paper controls for the generic effect of adding \emph{an} instruction to the
prompt. Third, the cost of skill text is framed as a budget problem, solved by progressive
disclosure or by moving skills into weights \citep{anthropic2025skills,yu2026latentskill}; the
hypothesis that domain-specific vocabulary actively interferes when a skill is used out of domain
has not been tested. To these we add a fourth: no paper audits its benchmarks for answer-structure
shortcuts that an optimiser can learn.

\subsection{The cost of instructions in context}

The motivating hypothesis for latent skills is \emph{context pollution}: a skill mined in one domain
brings that domain's vocabulary into another. Direct evidence for skill-level pollution does not
exist, but three adjacent findings make it plausible.

\paragraph{Irrelevant content degrades reasoning.} \citet{shi2023distracted} add irrelevant
sentences to grade-school arithmetic problems and find that accuracy drops substantially across
prompting techniques, and that an explicit instruction to ignore irrelevant information recovers
part of the loss. \citet{liu2024lost} show that models use information in long contexts unevenly,
performing best when relevant content is at the start or end. Both results concern distracting
\emph{content}; a cross-domain skill is distracting content with the added risk that it looks
relevant.

\paragraph{Performance is sensitive to features of the prompt that should not matter.}
\citet{sclar2024formatting} show that semantically equivalent changes in prompt formatting move
few-shot accuracy by large margins, and that the ranking of models can change with format. If
formatting alone moves results this much, then an instruction's presence, length and register are
confounds whenever a skill is compared against no skill.

\paragraph{Prompt meaning matters less than expected.} \citet{webson2022prompts} find that models
trained with prompts learn about as quickly from irrelevant or misleading instruction templates as
from instructive ones. This is the closest precedent in the prompting literature for the placebo
control used in our experiments: it suggests that part of what an instruction does is a consequence
of an instruction being there, independent of what it says. Our placebo results
(\S\ref{sec:shortcut}, \S\ref{sec:clean}) are a direct instance: a length-matched text with no
procedure matches the best skill on LiveMath OTHER, and on the shortcut-free suite placebos change
accuracy as much as skills do, mostly by changing reply length.

\subsection{Moving instructions into weights}

One response to the cost of context is to compile instructions into parameters. \emph{Context
distillation} \citep{snell2022distilling} fine-tunes a model so that its behaviour without a context
matches its behaviour with it. \emph{Gist tokens} \citep{mu2023gist} train a model to compress a
prompt into a few learned tokens that can be cached and reused. Parameter-efficient adapters such as
LoRA \citep{hu2022lora} make per-skill weight updates cheap, and hypernetworks can generate them from
text: \emph{Text-to-LoRA} \citep{charakorn2025t2l} produces a task-specific LoRA adapter from a
natural-language task description in a single forward pass.

\emph{LatentSkill} \citep{yu2026latentskill} applies this idea to agent skills and is the nearest
neighbour to our work. A Transformer hypernetwork maps each skill document to a set of LoRA updates,
trained first on a corpus of skill documents and then on agent trajectories. The resulting adapters
are modular and composable, remove skill text from the context (reducing context overhead
substantially), and improve over the same skills used in context on ALFWorld and seven Search-QA
datasets. Its in-context-skill ablation is the tier-1 versus tier-2 analogue of our tier-1 versus
tier-3 comparison. LatentSkill does not compare against activation steering, and it requires a
trained hypernetwork; an activation vector requires no training and is read off the model's own
response to the skill text.

\subsection{Activation steering}

\paragraph{Linear representations} Steering rests on the observation that many concepts are
represented approximately linearly in the residual stream. \citet{park2024linear} formalise the
\emph{linear representation hypothesis} and relate the directions used for probing to those used for
intervention. \citet{marks2024geometry} show emergent linear structure for truth and falsehood, and
argue that the difference of class means, which they call \emph{mass-mean probing}, identifies
directions that are more causally implicated in behaviour than those found by trained
logistic-regression probes, which can latch onto correlated features. This is the statistical
justification for extracting a skill by subtracting and averaging rather than by training a
classifier.

\paragraph{Extracting and applying steering vectors} \citet{subramani2022extracting} showed that a
vector \emph{optimised} to make a frozen model generate a target sentence can be added to hidden
states to steer generation, and that such vectors support arithmetic for style transfer.
\emph{Activation Addition} \citep{turner2023actadd} removed the optimisation: the difference between
activations on a single pair of prompts, added at inference, steers topic and sentiment. A single
pair is noisy, however, and conflates the target with whatever else differs between the prompts.

\emph{Inference-Time Intervention} \citep{li2023iti} identifies attention heads whose activations
linearly separate truthful from untruthful answers and shifts them along the separating direction at
inference, substantially improving TruthfulQA performance. \emph{Contrastive Activation Addition}
(CAA) \citep{panickssery2024caa} averages activation differences over a dataset of matched prompt
pairs at the answer-token position and adds the mean vector, scaled by a coefficient, at every
subsequent position; it steers sycophancy, hallucination, corrigibility and related behaviours in
Llama 2. \emph{Representation Engineering} \citep{zou2023repe} develops a family of reading and
control methods based on contrastive stimuli, including principal components of activation
differences. \emph{Mean-centring} \citep{jorgensen2023meancentring} subtracts the mean activation
over a generic dataset from the mean over a target dataset, an explicit baseline subtraction that
improves steering and is a methodological cousin of our placebo arm. \citet{arditi2024refusal} show
that refusal across many open models is mediated by a single difference-in-means direction, which
can be ablated to remove refusal or added to induce it; the paper is a practical precedent for layer
selection and coefficient sweeps.

Comparing these extraction choices, \citet{im2025steering} give a unified account of steering
methods and find that simple difference-of-means directions are competitive with, or better than,
PCA- and classifier-based alternatives. Later work conditions steering on the input:
\emph{conditional activation steering} \citep{lee2025cast} applies a behaviour vector only when the
input's activations match a condition vector. A separate line \emph{learns} the intervention: ReFT
\citep{wu2024reft} trains low-rank interventions on hidden representations as a parameter-efficient
alternative to fine-tuning, and RePS \citep{wu2025reps} trains steering vectors with a
reference-free preference objective. \emph{Persona vectors} \citep{chen2025persona} are extracted
automatically from a natural-language description of a trait, and are used both to monitor and to
steer that trait; this is a direct example of turning a text description into an activation
direction.

We use difference-in-means extraction in the CAA tradition, averaged over 50 prompts (40 on the
shortcut-free suite). Its injection is norm-relative,
$h_\ell \leftarrow h_\ell + \gamma\lVert h_\ell\rVert\, v/\lVert v\rVert$, so that the coefficient is
comparable across layers and only the vector's direction sets steering strength.

\paragraph{Tasks and instructions as vectors} Most steering work targets \emph{dispositions}: whether
to refuse, flatter, or be truthful. Closer to procedures, \emph{function vectors}
\citep{todd2024function} are the summed mean outputs of a small set of attention heads over
in-context examples of a task; added to a zero-shot prompt, they cause the model to execute that
task. \emph{Task vectors} \citep{hendel2023taskvectors} show that an in-context demonstration set
compresses into a single vector at one layer that can replace the demonstrations, and
\emph{in-context vectors} \citep{liu2024icv} use a latent vector derived from demonstrations to
replace them at inference. These results are the precedent for the equivalence our method relies on:
what text in context does can be summarised by a vector in the hidden state. They concern short
input--output mappings, however, not procedures that unfold over hundreds of generated tokens.

The most direct precedent for steering on instructions is \citet{stolfo2025instruction}, who compute
vectors as the difference between activations with and without an instruction and use them to impose
format, length and word-inclusion constraints, including in combination. This is the construction of
our ``vs nothing'' extraction arm. Our ``vs dummy'' arm adds what that design lacks: by subtracting a
length- and register-matched placebo system prompt with no procedural content, it separates the
effect of \emph{this} instruction from the effect of \emph{an} instruction.

\paragraph{Procedures and reasoning as vectors} Two recent lines move steering towards procedures.
\citet{venhoff2025reasoning} identify reasoning behaviours in thinking models, such as backtracking
and expressing uncertainty, and show that steering vectors can increase or suppress them. The
behaviours are procedural, but the evaluation is of the behaviours themselves rather than of task
accuracy in new domains.

\emph{Latent skill steering} \citep{jiang2026latent} is the closest prior work in method and the
only prior text-versus-vector comparison for procedural skills. It builds two controlled benchmarks:
SkillSet-Math, with 34 atomic and composite arithmetic and algebra skills, and SkillSet-Code, with
nine pairs of alternative algorithmic strategies such as depth-first versus breadth-first search.
Math-skill directions are the top principal component of last-token activations at every layer over
up to 500 examples of a skill, where each input contains the instruction, the problem and its answer
(for composite skills, the full derivation), a choice the authors disclose. Directions are
unit-normalised and added to the last half of layers at the last prefill token and every generated
token. On Qwen3-0.6B this improves accuracy on 29 of 34 skills, by 19.1 points on average. On 18
composite skills, an equal-weighted sum of the constituent atomic directions scores highest on 13,
averaging 77.1\% against 75.7\% for a concise text description and 74.2\% for a chain-of-thought
decomposition (51.9\% unsteered), on 100 test items per skill. Applied to GSM8K and MATH500
\citep{cobbe2021gsm8k,hendrycks2021math} across Qwen3 models from 0.6B to 14B, the directions match
or beat zero-shot in 8 of 10 cells, but the gains shrink with scale: on Qwen3-8B they are $+0.3$
points on GSM8K and $+0.8$ on MATH500. The paper reports no significance tests and no
random-direction or placebo control, and its transfer stays within mathematics. Our results at 8B to
27B are consistent with its shrinking gains: on the shortcut-free suite we find none.

\paragraph{Limits of steering} Two results temper the expectation that a vector will outperform the
prompt it was extracted from. \emph{AxBench} \citep{wu2025axbench} compares representation-based
steering with prompting and fine-tuning on concept steering at scale and finds prompting the
strongest method, ahead of every representation-based method evaluated, with sparse-autoencoder
features weakest; learned steering objectives such as RePS \citep{wu2025reps} narrow but do not close
this gap. \citet{tan2024reliability} find that CAA-style vectors are unreliable in ways averages
hide: steerability varies widely across inputs and sometimes reverses, part of a vector's apparent
effect reflects spurious biases such as a preference for a particular answer position, and
generalisation under distribution shift is mixed. A further, practical limit is degeneration: at
high coefficients, and especially when many layers are steered at once, generation collapses into
repetition \citep{turner2023actadd,panickssery2024caa}, which metrics computed on teacher-forced
tokens cannot detect.

\subsection{Methodological issues}

\paragraph{Controls for prompt presence.} CAA's matched pairs imply a control, since each pair
differs only in the target behaviour. Designs that contrast an instruction with no instruction
\citep{stolfo2025instruction} drop it, and the text-skill literature has none. Mean-centring
\citep{jorgensen2023meancentring} subtracts a generic baseline, but a generic corpus is not a matched
placebo instruction. The prompting evidence of \citet{webson2022prompts} and
\citet{sclar2024formatting} suggests such a control matters. We are not aware of prior work that
separates the effect of a particular instruction from that of an instruction being present; our
placebo-subtracted extraction arm does so.

\paragraph{Statistical comparison.} Skill effects on one benchmark are often a few points, close to
the standard error on a few hundred items. Paired comparisons on the same items with McNemar's test
\citep{mcnemar1947}, as recommended for comparing classifiers on one test set
\citep{dietterich1998}, and paired or clustered standard errors \citep{miller2024errorbars}, are the
appropriate tools; neither the skill-evolution papers nor latent skill steering report such tests.
Per-item fixed and broken counts and held-out subsets address the input-level variance documented by
\citet{tan2024reliability}.

\paragraph{Intrinsic versus task metrics.} Steering papers often choose layers and coefficients with
an intrinsic signal \citep{panickssery2024caa}. A KL-based fidelity score, which measures how
closely a steered model reproduces the text skill's token distribution, belongs to this family. It
is useful for choosing operating points but does not measure task value: if the text skill does not
improve accuracy, a faithful imitation will not either. Intrinsic and task metrics should be
reported separately.

\paragraph{Degeneration and truncation.} Teacher-forced metrics cannot see free generation
collapsing into repetition at high coefficients or across many layers
\citep{turner2023actadd,panickssery2024caa}, so an unforced diversity check is needed. A fixed
generation budget interacts with any intervention that changes verbosity; since skills change
verbosity, truncation introduces bias rather than noise.

\paragraph{Benchmark shortcuts.} An optimiser that scores edits on a benchmark will exploit whatever
regularity that benchmark has, including regularities in how its answers are constructed. On
LiveMath this produced a rule worth up to 72 points that contains no mathematics. Auditing a
benchmark with no model in the loop for answer-structure heuristics is a cheap precondition for
optimising skills against it.

\subsection{Benchmarks}

Our matrix benchmarks are drawn from the set used by WikiSkill \citep{tang2026wikiskill}:
LiveMathematicianBench \citep{he2026livemath}, SpreadsheetBench \citep{ma2024spreadsheetbench},
OfficeQA \citep{databricks2025officeqa,opsahlong2026officeqapro}, SealQA \citep{pham2025sealqa} and
ALFWorld \citep{shridhar2021alfworld}. They are also used, in various subsets, by EvoSkill,
Trace2Skill and SkillOpt, which makes our text-skill results directly comparable with that
literature. The shortcut-free suite (\S\ref{sec:suite}) adds MMLU-Pro, SuperGPQA, BeyondAIME,
HARDMath2, RealMath, IMO-AnswerBench and ArxivMath.

\subsection{Positioning}

Table~\ref{tab:positioning} places the closest work in the terms of our design. For each benchmark
$i$ and skill source $j$, a cell compares the text skill written for benchmark $j$ with the
embedding skill derived from it, evaluated on benchmark $i$. Diagonal cells ($i=j$) measure
in-domain effect; off-diagonal cells ($i\neq j$) measure transfer to a domain the skill was not
designed for.

No prior work reports the off-diagonal cells for activation-space skills. The contribution is
therefore not the extraction method, which is CAA, nor the in-domain text-versus-vector comparison,
which \citet{jiang2026latent} make on synthetic arithmetic. It is the comparison across unrelated
domains, at 8B to 27B scale, on the benchmarks the agent-skills literature uses, with a placebo
control for prompt presence and paired significance tests, and the finding that on benchmarks
without answer shortcuts neither representation improves on a length-matched placebo.

\begin{table*}[t]
\centering\small
\begin{tabular}{@{}p{3.4cm}p{2.3cm}p{3.2cm}p{2.8cm}p{2.5cm}@{}}
\toprule
Work & Representation & Cells of the skill-by-benchmark matrix & Domains & Presence control \\
\midrule
EvoSkill, Trace2Skill, SkillOpt, WikiSkill & Text & Diagonal; some off-diagonal & Agent-skills set & No \\
LatentSkill & Weights vs.\ text & Diagonal & ALFWorld, Search-QA & No \\
CAA, refusal direction, function vectors & Activations & Not applicable (no skills) & Behaviours, short tasks & Matched pairs (CAA) \\
Instruction steering & Activations & Diagonal & Format constraints & No \\
Latent skill steering & Activations vs.\ text & Diagonal, within mathematics & SkillSet, GSM8K, MATH500 & No \\
\textbf{Ours} & \textbf{Activations vs.\ text} & \textbf{Diagonal and transfer cells} & \textbf{Three domains + 7 math sets} & \textbf{Placebo arm and text placebo} \\
\bottomrule
\end{tabular}
\caption{Positioning against the closest prior work.}
\label{tab:positioning}
\end{table*}

\section{Method details}
\label{app:method}

\subsection{Development of the LiveMath text skill}
\label{app:livemathskill}

We wrote three versions of the LiveMath skill. Each changed behaviour measurably: the first drove
the rate of restating the claim before answering from 10\% to 100\% of replies, and the second drove
explicit relevance screening of the options from 0\% to 87\%. Accuracy did not move. LiveMath labels
the role of every distractor, and Table~\ref{tab:roles} shows why: each version closed one kind of
trap and opened another. The third version replaced two competing preferences with a single
criterion (is the option derivable from the hypotheses; strength ordering only after that), which
recovered the losses of the first two but did not exceed the baseline. It was written after seeing
the aggregate role histogram of versions one and two on the same 608 items, never per-item outcomes,
so it is test-set-informed; for this reason we report held-out results with every LiveMath number.

\begin{table}[t]
\centering\small
\begin{tabular}{@{}lrrr@{}}
\toprule
Answer chosen (\%) & Base & v1 & v2 \\
\midrule
Correct & 26.0 & 24.2 & 24.0 \\
Off-topic wildcard & 21.4 & 21.5 & 17.4 \\
Weaker but true & 18.3 & 11.8 & 20.7 \\
Stronger trap & 14.3 & 15.6 & 14.5 \\
Scope errors & 12.8 & 18.8 & 16.6 \\
\bottomrule
\end{tabular}
\caption{Distribution of chosen answers by distractor role on LiveMath (Qwen3-8B, $n=608$). Each
skill version redistributes errors across trap types without reducing them. v3 (26.0\% correct) is
omitted because its per-option breakdown was not saved.}
\label{tab:roles}
\end{table}

\subsection{Validity of the placebo}
\label{app:placebo}

Two checks support the placebo as a control for prompt presence. Behaviourally, the placebo used as a
text system prompt scores 26.15\% on LiveMath with Qwen3-8B, identical to the no-prompt baseline.
Geometrically, three placebos with unrelated content (brevity instructions; notation and typography;
tense, voice and naming) align with the skill direction to the same degree: cosine $+0.62$ on the 8B
and $+0.35$ to $+0.41$ on the 27B. If the alignment came from shared content it would vary with the
content; it does not, so the component shared by skill and placebo is the presence of a system
prompt.

A placebo is not neutral on every benchmark, however. On the shortcut-free suite the length-matched
placebos ask for brevity, halve reply length, and cost up to 9.4 points where problems need long
derivations; where the bare model stops early (HARDMath2, 27B) any system prompt lengthens replies
and the placebo gains 14.9 points (\S\ref{sec:clean}). A placebo therefore controls for prompt presence only where it does
not itself change verbosity, and its effect must be read next to reply length.

\subsection{Extraction pool}
\label{app:pool}

We compared an in-domain extraction pool with a domain-neutral one. The in-domain pool gave higher
fidelity in three of four model and arm combinations, and by a wide margin on the 27B ($+0.229$
against $+0.064$), so all matrix vectors are extracted from the benchmark's own items. On the diagonal
this means the pool items are also scored; the LiveMath 27B vs-nothing effect is $+6.6$ points on the
558 items outside the pool against $+6.25$ on all 608, so the overlap confers no advantage. On the
shortcut-free suite the pool is each benchmark's train split, so no pool item is scored.

\subsection{Degeneration guard}
\label{app:degen}

Table~\ref{tab:degen} lists the configurations that fidelity favoured but the unforced guard
rejected. The ceiling is $\gamma=0.3$ on both models; $\gamma=0.5$ is not usable, and $\gamma=1.0$,
the value in our original design, is roughly three times past it. Multi-layer spans are far more
sensitive than single layers. A first 27B sweep also placed every vector at layer 60, the deepest
layer searched; extending the grid to layer 63 moved the optimum to the interior, which is why we
check for boundary optima.

\begin{table}[t]
\centering\small
\begin{tabular}{@{}llrrr@{}}
\toprule
Model & Layer(s), $\gamma$ & Steered & Plain & Change \\
\midrule
8B & 34, 0.5 & 0.114 & 0.358 & $-68\%$ \\
27B & 62, 0.5 & 0.273 & 0.413 & $-34\%$ \\
27B & 60--63, 0.2 & 0.008 & 0.407 & $-98\%$ \\
\bottomrule
\end{tabular}
\caption{Operating points rejected by the degeneration guard: distinct-token ratio of unforced
generations with and without steering.}
\label{tab:degen}
\end{table}

\subsection{Budgets and answer forcing}
\label{app:budget}

Token budgets differ by model because larger models deliberate longer. At a 2048-token budget, 51.6\%
of Qwen3.6-27B replies on LiveMath reached the cap before stating an answer, and the model scored
16.94\%, below the 20\% chance rate and below the 8B. With a 4096-token budget and answer forcing its
baseline is 33.06\%. Forcing is applied to 2.6\% of 8B replies and 21.7\% of 27B replies on
LiveMath; forced rates for every cell are reported with the results. The SpreadsheetBench agent
variant runs for up to 20 turns with no per-turn cap. On the shortcut-free suite at 4096 tokens,
17--25\% of BeyondAIME replies never reach a final answer, and the rate differs by condition; forcing
is applied there as well.

\subsection{Harness faults}
\label{app:faults}

We report the faults in our own harness because each would have reversed a conclusion without
producing an error.

\begin{enumerate}
  \item \textbf{Truncation.} Scoring capped replies as wrong made the 27B look worse than the 8B
  (\S\ref{app:budget}).
  \item \textbf{Tool-call parser.} The agent server parsed tool calls in the XML format of Qwen3.6,
  while Qwen3-8B and Qwen3-14B emit JSON inside \texttt{<tool\_call>} tags. On this mismatch the
  server returns no tool calls, the agent executes nothing, and the benchmark scores exactly 0\% with
  nothing logged. It was found by querying the live server directly; the parser must be set per model
  generation.
  \item \textbf{Silent shard failure.} The evaluator ran four GPU shards and waited on them jointly,
  which returned only the last shard's status. One shard failed on a transient filesystem error and
  the job reported success with a quarter of the results missing, detected only because the report
  printed $n=456$ rather than 608. Each shard's status is now checked.
  \item \textbf{Filename collision.} Output files carried no model tag, so a run of one model
  overwrote the results of the other.
  \item \textbf{Math-Verify not installed.} The environment that ran the shortcut-free suite lacked
  the \texttt{math\_verify} package. The scorer caught the import error as a failed check, so every
  free-form cell counted exact string matches only (no symbolic or numeric match in 16.7k test
  records). All free-form cells are rescored from the saved replies with it installed; accuracies rise
  by up to 19 points (HARDMath2 27B, no skill: 4.1\% to 23.1\%), and on HARDMath2 27B the rescoring
  reveals effects of 8 to 15 points that string matching had hidden.
  \item \textbf{Math-Verify in threads.} Math-Verify's timeouts use \texttt{signal.alarm}, which
  fails silently outside the main thread; inside SkillOpt's rollout threads every symbolic and
  numeric check returned ``wrong'', and only exact string matches counted. Comparisons now run in a
  process pool under our own deadline.
  \item \textbf{Gold-answer formatting.} RealMath writes answers inside \verb|\( \)| and \verb|\[ \]|
  and as equations (\verb|\omega = 1| for \verb|1|); a 6-digit decimal was accepted for an exact
  rational under the default tolerance. Delimiters are stripped, right-hand sides compared, and
  rounding tightened to 10 digits.
  \item \textbf{Answer-format clash.} SkillOpt's skills ask for \texttt{<answer>X</answer>} while the
  evaluation prompt asks for \texttt{Answer: <letter>}; models wrote hybrids that a strict parser
  scored as unanswered, costing optimised drafts up to 30 points on dev for a format clash. The
  parser now reads every committed form, and all results were re-scored from saved replies.
  \item \textbf{Silent garbage from a GPU.} After a reboot one GPU produced only repeated tokens
  (64 of 64 replies) with no error; two SkillOpt steps were rolled back and redone, and a two-minute
  generation probe now runs before any job.
  \item \textbf{Memory in the forcing pass, logged as success.} On SuperGPQA with the 27B the forcing
  pass (prompt plus a 4096-token reply) ran out of memory because the retry ran while the exception
  still held the failed attempt's tensors, and a logging bug recorded exit code 0. The cells were
  rerun after both were fixed.
\end{enumerate}

\subsection{SkillOpt runs on LiveMath without the meta option}
\label{app:skillopt}

\begin{table*}[t]
\centering\small
\begin{tabular}{@{}llrrrrrr@{}}
\toprule
SkillOpt skill & Model & Base & Placebo & \TS & \ES{} vs-placebo & \ES{} vs-nothing & One-token vector \\
\midrule
8B-trained, with meta rule & 4B & 41.7 & +5.6 & +0.0 & +7.1$^*$ & +4.3 & +1.5 \\
 & 8B & 43.5 & +1.9 & $-$1.9 & $-$2.5 & $-$0.6 & +0.0 \\
 & 14B & 43.5 & +4.6 & $-$1.2 & $-$0.9 & +0.3 & +1.2 \\
 & 27B & 52.2 & +1.2 & $-$2.2 & $-$2.8 & +1.9 & $-$4.0 \\
\midrule
8B-trained, meta rule removed & 4B & 41.7 & +9.0$^*$ & +1.2 & $-$1.2 & +4.0 & +2.2 \\
 & 8B & 43.5 & +3.4 & +2.2 & $-$6.8$^*$ & $-$1.2 & +1.5 \\
 & 14B & 43.5 & +6.8$^*$ & +1.9 & $-$8.3$^*$ & +1.5 & $-$1.2 \\
 & 27B & 52.2 & +0.6 & +5.9$^*$ & $-$1.5 & $-$0.9 & $-$2.2 \\
\midrule
27B-trained, META items excluded & 4B & 41.7 & +0.9 & +6.2$^*$ & +3.1 & $-$0.3 & +2.5 \\
 & 8B & 43.5 & +0.0 & +4.3 & +0.0 & $-$2.8 & +0.9 \\
 & 14B & 43.5 & +6.2$^*$ & +0.0 & +1.2 & +1.2 & $-$2.8 \\
 & 27B & 52.2 & +0.3 & +0.0 & $-$0.9 & $-$0.3 & $-$3.1 \\
\midrule
27B-trained, meta edits filtered & 4B & 41.7 & +0.0 & +2.2 & $-$0.6 & $-$3.1 & +2.2 \\
 & 8B & 43.5 & +2.2 & $-$0.6 & $-$2.5 & $-$2.2 & $-$1.2 \\
 & 14B & 43.5 & +2.8 & +4.0 & $-$0.3 & +1.5 & $-$0.3 \\
 & 27B & 52.2 & $-$1.9 & +2.2 & +0.9 & $-$1.5 & $-$2.5 \\
\midrule
27B-trained, v2 (partial, $n=153$) & 27B & 51.2 & +7.8$^*$ & $-$2.6 & $-$3.9 & $-$4.9$^\dagger$ & $-$3.9$^\dagger$ \\
\bottomrule
\end{tabular}
\caption{SkillOpt skills learned on LiveMath without relying on the meta option: change in accuracy
(points) on the OTHER subset ($n=324$) against the no-skill baseline. Placebo is length-matched text
with no procedure; the one-token vector compresses the whole skill into a single steering direction.
$^{*}$ $p<0.05$, exact McNemar. $^\dagger$ $n=103$.}
\label{tab:skillopt}
\end{table*}
Table~\ref{tab:skillopt} gives the OTHER-subset results for each variant. The
8B-trained skill with its rule takes META to 68.5\%, 94.4\%, 99.1\% and 97.4\% on the 4B, 8B, 14B
and 27B while OTHER changes by $-1.9$ to $0.0$. SkillOpt's own split makes the shortcut hard to
avoid: 31 of its 52 train and validation items are META.

\subsection{Audit of the learned skills}
\label{app:audit}

\begin{table*}[t]
\centering\small
\begin{tabular}{@{}p{0.2\textwidth}p{0.45\textwidth}p{0.28\textwidth}@{}}
\toprule
Skill & Content that exploits the benchmark & Effect \\
\midrule
\texttt{strength\_audit} & Yes: a rule for the meta option & Its whole LiveMath gain \\
\texttt{livemath\_ts} & Partly: notes that wrong options usually differ from the right one in one
detail (the basis of the centroid), never says to use it & None on LiveMath beyond the placebo; on
free-form questions it makes small models invent options \\
SkillOpt, LiveMath (GPT-5.5 and ours) & Yes: ``meta-option absolute priority'' & The LiveMath gain \\
SkillOpt, MCQ suite & Mostly method; two memorised training items (an answer key for $\sqrt[3]{130}$,
Fermat primes for a polygon question) and one guessing rule (``pick the option with the same
structure as your result'') & No gain on centroid-right items; neither memorising version chosen \\
SkillOpt, free-form suite & Method plus benchmark-construction priors (HARDMath2: ``reject constant
answers'', ``amplitude is always $Q^{-1/4}$''; RealMath: ``never answer `cannot be determined'\,''),
anti-loop rules & None significant; not chosen \\
Placebos & No procedure; ``be concise'' style rules & Shorter replies; losses where problems need room \\
\bottomrule
\end{tabular}
\caption{What the evaluated skills contain, read line by line.}
\label{tab:audit}
\end{table*}

\subsection{Compute}

The matrix runs used single H100 GPUs, with four independent shards per cell. The shortcut-free
suite and its vector arm ran on two NVIDIA RTX PRO 6000 Blackwell GPUs (96\,GB), with vLLM 0.32 for
text conditions and HuggingFace \texttt{transformers} for every vector-arm cell; each 27B vector-arm
cell took about 2.5 hours. 

\end{document}
